% !TEX encoding = UTF-8 Unicode

\documentclass[twocolumn,10pt,a4j]{ltjsarticle}
\usepackage{luatexja}
\usepackage{luatexja-fontspec}
\usepackage{kougai}

\pagestyle{empty}

\title{目的別スクワットフォーム修正システムの開発}
\author{2332178 渡辺　大地　　指導教員 須田 宇宙 准教授}
\date{}

\begin{document}

\maketitle

\section{はじめに}
筋力トレーニングは怪我の予防や健康維持に効果的な手段として広く認識されている．
シニア世代の運動目的では健康維持が9割を占め，運動により筋力を維持したい場所は脚90.1\%，体幹(腰を含む)65.4\%であった\cite{senior}．
脚，体幹を同時に鍛える筋力トレーニング方法としてスクワットがある．
スクワットは，立位から膝と股関節を屈伸させる動作を繰り返すレジスタンス運動である．
スクワットは複数の関節を同時に動かす多関節種目であり，フォームの安定しない初心者や高齢者の場合一つの関節に負荷が集中してしまい怪我の原因になる．
また，スクワットでは，負荷を軽減させたい関節や優先的に働かせたい筋肉によって最適なフォームが異なる．

現在，スクワットのフォームを修正するための支援システムはすでに存在する．
しかし，それらの多くは一般的なフォームを基準として修正を行うものであり，利用者の目的に適したフォームを提示するシステムは確認されていない．

そこで本研究では，カメラ映像から算出した骨格情報を元に，適切なフォームへと修正するためのフィードバックを行うシステムの開発を目的とする．

\section{スクワットについて}
スクワットでは，足幅を肩幅程度にとり，つま先をやや外側に向け，背筋を伸ばした姿勢ではじめる．
膝と股関節を同時に曲げ，上体を自然な範囲で前傾させ，足全体に重心がかかるようにしながら太ももが床と平行になる程度まで屈伸するのが一般的なフォームである．
スクワット動作においては，足幅，体幹の前傾角度，すねの角度，つま先の向き，重心の位置，屈伸する深さ，さらには股関節主導か膝関節主導かといった要素によって，膝関節や腰関節にかかる負荷や，働く筋肉の割合が変化する\cite{squat}．

膝関節への負担軽減を目的とする場合，足幅を広めにとり，体幹を前傾させ，すねの前傾を抑えるフォームとする．

深く屈伸することで股関節の最大可動域に達し，それ以上の深さで屈伸を行うと腰を曲げて体幹を前傾させてしまうため，腰への負担が増加する．
腰への負担軽減を目的とする場合，足幅をやや広げ，体幹をできるだけ垂直に保ち，すねを前傾させ，深く屈伸しないフォームとする．
また，屈伸が深くなるほど負荷は全体的に増加するため，膝関節への負担軽減を目的とする場合でも，膝に痛みがある場合は深く屈伸しないフォームにする必要がある．

\begin{figure}[h]
 \begin{center}
  \includegraphics[clip,width=80mm,height=50mm]{motion_capture_next.pdf}
 \end{center}
 \vspace{-1.5em}
 \caption{骨格情報を正面(左)と横側(右)から取得した画像}
 \label{fig:MediaPipe}
 \vspace{-2em}
\end{figure}

\section{スクワットの判定}
現在，スクワットのフォームを修正するための支援システムはすでに存在するが，それらの多くは一般的なフォームを基準として修正を行うものである．
そのため腰保護や膝関節保護などに適したフォームを提示するシステムは確認されていない．

本研究では，膝関節や腰関節への負担軽減を目的としたフォームを提示するために，既存のシステムに加えて足幅，肩幅，つま先の向きの要素を取得する．

図\ref{fig:MediaPipe}はカメラ映像から取得した骨格情報を正面と横側から表示したものである．
既存のシステムでは，図\ref{fig:MediaPipe}(右)のように，横側からのカメラ映像で膝関節の角度，股関節の角度，上体の角度，太ももの角度を取得し，それらの数値が許容範囲内であった場合にスクワットと判定している．
そのため，それらの要素は図\ref{fig:MediaPipe}(左)のように，正面からのカメラ映像を用いる．

%スクワットを判定するには，足幅，肩幅，つま先の向き，膝関節の角度，股関節の角度，上体の角度，太ももの角度の要素が必要になる．
%図\ref{fig:MediaPipe}のように足幅，肩幅の要素は正面からのカメラ映像から取得できる．
%また，膝関節の角度，股関節の角度，上体の角度，太ももの角度は横側からのカメラ映像から取得できる．
%つま先の向きは，正面を向いている状態を初期状態としたときの関節角度の変化で取得することができる．
%各要素の数値が許容範囲内であった場合にスクワットと判定する．

関節座標の取得は，オープンソースの機械学習フレームワークである MediaPipe を用いて，カメラ映像から人体の関節座標をリアルタイムに取得する．

\section{今後の予定}
各目的に適したフォームを各要素ごとに定め，スクワット動作中の関節座標を取得するシステムを構築する．

%\pagebreak
\begin{thebibliography}{99}
\bibitem{senior} 生き活きライフ(フレイル対策)部会: ``＜シニアの運動に関する意識調査＞毎日運動していても9割以上が筋力の低下を実感``,
\url{https://prtimes.jp/main/html/rd/p/000000001.000146886.html}， 2026/7/8参照

\bibitem{squat} Straub RK,Powers CM: ``A Biomechanical Review of the Squat Exercise: Implications for Clinical Practice.'', International Journal of Sports Physical Therapy. 2024;19(4):490-501，
\url{https://doi.org/10.26603/001c.94600}，2026/7/15参照

\end{thebibliography}
 
\end{document}

